\subsection{广义函数}

我们沿用上一节的记号；于是 U 表示 \(R^{n}\) 的开子集。

定义 1. --- \(\mathcal{D}(\mathrm{U})\) 的对偶空间，赋以有界收敛拓扑，称为 U 上的广义函数空间。记作 \(\mathcal{D}'(\mathrm{U})\) 。

于是，U 上的广义函数就是 \(\mathcal{D}(\mathrm{U})\) 上的连续线性形式。

若 \(f\) 是 \(\mathcal{D}(\mathrm{U})\) 上的线性形式（特别地，若 \(f\) 是广义函数），且 \(\varphi\) 属于 \(\mathcal{D}(\mathrm{U})\) ，我们用 \(\langle f, \varphi \rangle\) 表示 \(f\) 在 \(\varphi\) 处的取值。

由于 \(\mathcal{D}(\mathrm{U})\) 是有界型的，空间 \(\mathcal{D}'(\mathrm{U})\) 是完备的（TVS, III, p. 24，推论 1）。由于 \(\mathcal{D}(\mathrm{U})\) 是蒙泰尔空间，\(\mathcal{D}'(\mathrm{U})\) 亦然（EVT, IV, p. 19，命题 9）。特别地，空间 \(\mathcal{D}'(\mathrm{U})\) 是自反的（EVT, IV, p. 16，定理 2）。

引理 4. --- 设 f 是从 \(\mathcal{D}(\mathrm{U})\) 到 \(\mathbf{C}\) 中的线性映射。那么 \(f\) 是广义函数当且仅当：对每个紧子集 \(\mathrm{K}\) （\(\mathrm{U}\) 的）以及每个族 \((\mathrm{M}_{\alpha})_{\alpha \in \mathbf{N}^{n}}\) （\(\mathbf{R}_{+}\) 中的），线性形式 \(f\) 在由这样的函数 \(\varphi \in \mathcal{C}_{\mathrm{K}}^{\infty}(\mathrm{U})\) ——对每个 \(\alpha \in \mathbf{N}^{n}\) 有

\[
\sup _ {x \in \mathrm{K}} | \partial^ {\alpha} \varphi (x) | \leqslant \mathrm{M} _ {\alpha}.
\]

——所成之集上有界。事实上，空间 \(\mathcal{D}(\mathrm{U})\) 是有界型的，且 \(\mathcal{D}(\mathrm{U})\) 的每个有界子集都含于陈述中所描述的有界集之一。

引理 5. --- 设 F 是 \(\mathcal{D}'(\mathrm{U})\) 上具有可数基或包含一个单有界集合的滤子。那么 F 收敛于某个广义函数当且仅当 \(\langle f, \varphi \rangle\) 对 U 上的每个测试函数 \(\varphi\) 依 F 收敛。

特别地，广义函数序列 \((f_{m})_{m\in\mathbf{N}}\) 收敛于广义函数 f 当且仅当：对每个 \(\varphi\in\mathcal{D}(\mathrm{U})\) ，序列 \((\langle f_{m},\varphi\rangle)_{m\in\mathbf{N}}\) 收敛于 \(\langle f,\varphi\rangle\) 。

由于 \(\mathcal{D}'(\mathrm{U})\) 是蒙泰尔空间，故为桶型空间；这由巴拿赫–斯坦豪斯定理得出（EVT, III, p. 26，定理 1 的推论 3）。

设 V 是含于 U 的开集。从 \(\mathcal{D}(\mathrm{V})\) 到 \(\mathcal{D}(\mathrm{U})\) 中的典范线性映射的转置是从 \(\mathcal{D}'(\mathrm{U})\) 到 \(\mathcal{D}'(\mathrm{V})\) 中的连续线性映射，称为广义函数从 U 到 V 的限制，记作 \(r_{\mathrm{VU}}\) ，有时也记作 \(r_{\mathrm{V,U}}\) 。

有 \(r_{\mathrm{VV}} = 1_{\mathcal{D}'(\mathrm{V})}\) 。若 \(W \subset V \subset U\) 是 U 的开子集，则我们有 \(r_{\mathrm{WV}} \circ r_{\mathrm{VU}} = r_{\mathrm{WU}}\) 。换言之，若用 \(\mathcal{T}\) 表示 U 的拓扑，则投射系 \(\underline{\mathcal{D}'}(\mathrm{U}) = ((\mathcal{D}'(\mathrm{V}))_{\mathrm{V} \in \mathcal{T}}, (r_{\mathrm{WV}}))\) 是 U 上取值于局部凸拓扑向量空间的结构种的预层（TA, I, p. 42，定义 1 及 p. 66, §10）。

命题 5. --- 预层 \(\underline{\mathcal{D}}^{\prime}(\mathrm{U})\) 是一个层。

回顾（TA, I, p. 43，定义 2），这意味着对 U 的每个开子集 V 及 V 的每个开覆盖 \((\mathrm{V}_{i})_{i\in\mathrm{I}}\) ，下列条件满足：

\begin{enumerate}
\def\labelenumi{(\roman{enumi})}
\item
  映射 \((r_{\mathrm{V}_i\mathrm{V}})_{i\in \mathrm{I}}\colon \mathcal{D}'(\mathrm{V})\to \prod_{i\in \mathrm{I}}\mathcal{D}'(\mathrm{V}_i)\) 是单射；
\item
  对每个族 \((f_i)\in \prod_{i\in \mathrm{I}}\mathcal{D}'(\mathrm{V}_i)\) ——满足
\end{enumerate}

\[
r _ {(\mathrm{V} _ {i} \cap \mathrm{V} _ {i ^ {\prime}}) \mathrm{V} _ {i}} (f _ {i}) = r _ {(\mathrm{V} _ {i} \cap \mathrm{V} _ {i ^ {\prime}}) \mathrm{V} _ {i ^ {\prime}}} (f _ {i ^ {\prime}})
\]

——对每对 \((i, i') \in \mathrm{I} \times \mathrm{I}\) 都成立，则存在广义函数 \(f \in \mathcal{D}'(\mathrm{V})\) ，使得对每个 \(i \in I\) ，有 \(r_{\mathrm{V}_{i}\mathrm{V}}(f) = f_{i}\) 。

设 V 是 U 的开子集，\(\mathcal{V} = (\mathrm{V}_{i})_{i \in \mathrm{I}}\) 是 V 的开覆盖。设 \((\mathrm{W}_{j})_{j \in \mathrm{J}}\) 是 V 的开覆盖，它局部有限、比 \(\mathcal{V}\) 更细，且由相对紧的部分组成（IV, p. 201，引理 2）。取定单位分解 \((\varphi_{j})_{j \in \mathrm{J}}\) ——从属于覆盖 \((\mathrm{W}_{j})_{j \in \mathrm{J}}\) ——使得 \(\varphi_{j}\) 的支集含于 \(W_{j}\) （对所有 \(j \in J\) ；IV, p. 196，引理 1）。

设 \(\varphi\in\mathcal{D}(V)\) 。由于这样的 \(j\in J\) ——使 \(W_{j}\) 与 \(\varphi\) 的支集相交——组成的集合是有限的（TG, I, §9, no. 1, p. 59），我们有

\[
\varphi = \sum_ {j \in \mathrm{J}} \varphi \varphi_ {j}
\]

其中和式只有有限多个非零项。

我们来证明 (i)。设 \(f \in \mathcal{D}'(\mathrm{V})\) 满足 \(r_{\mathrm{V}_i\mathrm{V}}(f) = 0\) （对所有 \(i \in \mathrm{I}\) ）。这意味着 \(\langle f, \varphi \rangle = 0\) 对每个测试函数 \(\varphi\) （其支集含于某个开集 \(\mathrm{V}_i\) ）成立。若 \(\varphi\) 的支集含于某个开集 \(\mathrm{W}_j\) ，则更是如此。但这时，对每个 \(\varphi \in \mathcal{D}(\mathrm{V})\) ，我们有

\[
\langle f, \varphi \rangle = \langle f, \sum_ {j \in \mathrm{J}} \varphi \varphi_ {j} \rangle = \sum_ {j \in \mathrm{J}} \langle f, \varphi \varphi_ {j} \rangle = 0,
\]

因而 f = 0。

我们来证明 (ii)。设 \((f_{i})_{i\in\mathrm{I}}\) 是这样的族：\(f_{i}\in\mathcal{D}^{\prime}(\mathrm{V}_{i})\) （对所有 \(i\in I\) ），且 \(r_{\mathrm{V}_{i}\cap\mathrm{V}_{i^{\prime}},\mathrm{V}_{i}}(f_{i})=r_{\mathrm{V}_{i}\cap\mathrm{V}_{i^{\prime}},\mathrm{V}_{i^{\prime}}}(f_{i^{\prime}})\) 对 I 中所有 i 与 \(i^{\prime}\) 成立。设 \(\iota:J\to I\) 是满足 \(\mathrm{W}_{j}\subset\mathrm{V}_{\iota(j)}\) （对所有 \(j\in J\) ）的映射。对每个 \(j\in J\) ，令 \(\widetilde{f}_{j}=r_{\mathrm{W}_{j},\mathrm{V}_{\iota(j)}}(f_{\iota(j)})\) 。我们于是有

\[
\begin{array}{r l} & r _ {\mathrm{W} _ {j} \cap \mathrm{W} _ {j ^ {\prime}}, \mathrm{W} _ {j}} (\widetilde {f} _ {j}) = r _ {\mathrm{W} _ {j} \cap \mathrm{W} _ {j ^ {\prime}}, \mathrm{W} _ {j}} (r _ {\mathrm{W} _ {j}, \mathrm{V} _ {\iota (j)}} (f _ {\iota (j)})) \\ & \qquad = r _ {\mathrm{W} _ {j} \cap \mathrm{W} _ {j ^ {\prime}}, \mathrm{V} _ {\iota (j)}} (f _ {\iota (j)}) \\ & \qquad = r _ {\mathrm{W} _ {j} \cap \mathrm{W} _ {j ^ {\prime}}, \mathrm{V} _ {\iota (j)} \cap \mathrm{V} _ {\iota (j ^ {\prime})}} \circ r _ {\mathrm{V} _ {\iota (j)} \cap \mathrm{V} _ {\iota (j ^ {\prime})}, \mathrm{V} _ {\iota (j)}} (f _ {\iota (j)}). \end{array}
\]

交换 j 与 j' 的角色，并注意到

\[
r _ {\mathrm{V} _ {\iota (j)} \cap \mathrm{V} _ {\iota (j ^ {\prime})}, \mathrm{V} _ {\iota (j)}} (f _ {\iota (j)}) = r _ {\mathrm{V} _ {\iota (j)} \cap \mathrm{V} _ {\iota (j ^ {\prime})}, \mathrm{V} _ {\iota (j ^ {\prime})}} (f _ {\iota (j ^ {\prime})}),
\]

由假设成立，我们推出

\[
r _ {\mathrm{W} _ {j} \cap \mathrm{W} _ {j ^ {\prime}}, \mathrm{W} _ {j}} (\widetilde {f} _ {j}) = r _ {\mathrm{W} _ {j} \cap \mathrm{W} _ {j ^ {\prime}}, \mathrm{W} _ {j}} (\widetilde {f} _ {j ^ {\prime}})\tag{1}
\]

对所有 \((j, j') \in \mathrm{J}^{2}\) 成立。

对 \(\varphi \in \mathcal{D}(\mathrm{V})\) ，令

\[
\lambda (\varphi) = \sum_ {j \in \mathrm{J}} \langle \widetilde {f} _ {j}, \varphi \varphi_ {j} | \mathrm{W} _ {j} \rangle ,
\]

其中和式是有限的，因为只有有限多项可能非零。

映射 \(\lambda\) 是 \(\mathcal{D}(\mathrm{V})\) 上的线性形式。我们来证明它是广义函数。设 B 是 \(\mathcal{D}(\mathrm{V})\) 的有界子集，K 是 V 的紧子集且 \(B \subset \mathcal{C}_{K}^{\infty}(V)\) 。依据 TG, I, §9, no. 1, p. 59，存在 J 的有限子集 J' 使得

\[
\lambda (\varphi) = \sum_ {j \in J ^ {\prime}} \left\langle \widetilde {f} _ {j}, \varphi \varphi_ {j} \mid W _ {j} \right\rangle
\]

对所有 \(\varphi\in\mathrm{B}\) 成立。由于 \(\widetilde{f}_{j}\) 对每个 \(j\in\mathrm{J}^{\prime}\) 是广义函数，且 \(\mathcal{D}(\mathrm{V})\) 是拓扑代数，可知 \(\lambda\) 在 B 上有界，由此得结论（引理 4）。

设 \(j \in J\) ，\(\varphi \in \mathcal{D}(V)\) 的支集含于 \(W_{j}\) 。我们于是有

\[
\langle \lambda , \varphi \rangle = \sum_ {j ^ {\prime} \in \mathrm{J}} \left\langle \widetilde {f} _ {j ^ {\prime}}, \varphi \varphi_ {j ^ {\prime}} \right| \mathrm{W} _ {j} \rangle = \sum_ {j ^ {\prime} \in \mathrm{J}} \left\langle \widetilde {f} _ {j}, \varphi \varphi_ {j ^ {\prime}} \right| \mathrm{W} _ {j} \rangle
\]

这依据 (1)，因为 \(\varphi\varphi_{j'}\) 的支集含于 \(W_j \cap W_{j'}\) 。从而

\[
\langle \lambda , \varphi \rangle = \langle \widetilde {f} _ {j}, \sum_ {j ^ {\prime} \in \mathrm{J}} \varphi \varphi_ {j ^ {\prime}} | \mathrm{W} _ {j} \rangle = \langle \widetilde {f} _ {j}, \varphi | \mathrm{W} _ {j} \rangle ,
\]

因而 \(r_{\mathrm{W}_j\mathrm{V}}(\lambda) = \widetilde{f}_j\) 对所有 \(j\in \mathbf{J}\) 成立。

设 \(i \in I\) 。最后我们来证明 \(\lambda\) 在 \(V_i\) 上的限制与 \(f_i\) 一致。将条件 (i) 应用于 \(V_i\) 的由诸开集 \(W_j\) 组成的覆盖，只需验证对每个 \(j \in J\) ，\(\lambda\) 在 \(V_i \cap W_j\) 上的限制与 \(f_i\) 的限制一致。由前所述，只需验证 \(f_i\) 在 \(V_i \cap W_j\) 上的限制与 \(\widetilde{f}_j\) 的限制一致。而我们有

\[
\begin{array}{c} r _ {\mathrm{V} _ {i} \cap \mathrm{W} _ {j}, \mathrm{V} _ {i}} (f _ {i}) = r _ {\mathrm{V} _ {i} \cap \mathrm{W} _ {j}, \mathrm{V} _ {i} \cap \mathrm{V} _ {\iota (j)}} (r _ {\mathrm{V} _ {i} \cap \mathrm{V} _ {\iota (j)}, \mathrm{V} _ {i}} (f _ {i})) \\ = r _ {\mathrm{V} _ {i} \cap \mathrm{W} _ {j}, \mathrm{V} _ {i} \cap \mathrm{V} _ {\iota (j)}} (r _ {\mathrm{V} _ {i} \cap \mathrm{V} _ {\iota (j)}, \mathrm{V} _ {\iota (j)}} (f _ {\iota (j)})) = r _ {\mathrm{V} _ {i} \cap \mathrm{W} _ {j}, \mathrm{V} _ {\iota (j)}} (f _ {\iota (j)}), \end{array}
\]

其中我们用到了关于族 \((f_{i})\) 的假设。于是

\[
r _ {\mathrm{V} _ {i} \cap \mathrm{W} _ {j}, \mathrm{V} _ {\iota (j)}} (f _ {\iota (j)}) = r _ {\mathrm{V} _ {i} \cap \mathrm{W} _ {j}, \mathrm{W} _ {j}} (r _ {\mathrm{W} _ {j}, \mathrm{V} _ {\iota (j)}} (f _ {\iota (j)})) = r _ {\mathrm{V} _ {i} \cap \mathrm{W} _ {j}, \mathrm{W} _ {j}} (\widetilde {f} _ {j}),
\]

由此即可得出结论。

